\IfFileExists{stacks-project.cls}{%
\documentclass{stacks-project}
}{%
\documentclass{amsart}
}

% For dealing with references we use the comment environment
\usepackage{verbatim}
\newenvironment{reference}{\comment}{\endcomment}
%\newenvironment{reference}{}{}
\newenvironment{slogan}{\comment}{\endcomment}
\newenvironment{history}{\comment}{\endcomment}

% For commutative diagrams we use Xy-pic
\usepackage[all]{xy}

% We use 2cell for 2-commutative diagrams.
\xyoption{2cell}
\UseAllTwocells

% We use multicol for the list of chapters between chapters
\usepackage{multicol}

% This is generally recommended for better output
\usepackage{lmodern}
\usepackage[T1]{fontenc}

% For cross-file-references
\usepackage{xr-hyper}

% Package for hypertext links:
\usepackage{hyperref}

% For any local file, say "hello.tex" you want to link to please
% use \externaldocument[hello-]{hello}
\externaldocument[introduction-]{introduction}
\externaldocument[conventions-]{conventions}
\externaldocument[sets-]{sets}
\externaldocument[categories-]{categories}
\externaldocument[topology-]{topology}
\externaldocument[sheaves-]{sheaves}
\externaldocument[sites-]{sites}
\externaldocument[stacks-]{stacks}
\externaldocument[fields-]{fields}
\externaldocument[algebra-]{algebra}
\externaldocument[brauer-]{brauer}
\externaldocument[homology-]{homology}
\externaldocument[derived-]{derived}
\externaldocument[simplicial-]{simplicial}
\externaldocument[more-algebra-]{more-algebra}
\externaldocument[smoothing-]{smoothing}
\externaldocument[modules-]{modules}
\externaldocument[sites-modules-]{sites-modules}
\externaldocument[injectives-]{injectives}
\externaldocument[cohomology-]{cohomology}
\externaldocument[sites-cohomology-]{sites-cohomology}
\externaldocument[dga-]{dga}
\externaldocument[dpa-]{dpa}
\externaldocument[sdga-]{sdga}
\externaldocument[hypercovering-]{hypercovering}
\externaldocument[schemes-]{schemes}
\externaldocument[constructions-]{constructions}
\externaldocument[properties-]{properties}
\externaldocument[morphisms-]{morphisms}
\externaldocument[coherent-]{coherent}
\externaldocument[divisors-]{divisors}
\externaldocument[limits-]{limits}
\externaldocument[varieties-]{varieties}
\externaldocument[topologies-]{topologies}
\externaldocument[descent-]{descent}
\externaldocument[perfect-]{perfect}
\externaldocument[more-morphisms-]{more-morphisms}
\externaldocument[flat-]{flat}
\externaldocument[groupoids-]{groupoids}
\externaldocument[more-groupoids-]{more-groupoids}
\externaldocument[etale-]{etale}
\externaldocument[chow-]{chow}
\externaldocument[intersection-]{intersection}
\externaldocument[pic-]{pic}
\externaldocument[weil-]{weil}
\externaldocument[adequate-]{adequate}
\externaldocument[dualizing-]{dualizing}
\externaldocument[duality-]{duality}
\externaldocument[discriminant-]{discriminant}
\externaldocument[derham-]{derham}
\externaldocument[local-cohomology-]{local-cohomology}
\externaldocument[algebraization-]{algebraization}
\externaldocument[curves-]{curves}
\externaldocument[resolve-]{resolve}
\externaldocument[models-]{models}
\externaldocument[functors-]{functors}
\externaldocument[equiv-]{equiv}
\externaldocument[pione-]{pione}
\externaldocument[etale-cohomology-]{etale-cohomology}
\externaldocument[proetale-]{proetale}
\externaldocument[relative-cycles-]{relative-cycles}
\externaldocument[more-etale-]{more-etale}
\externaldocument[trace-]{trace}
\externaldocument[crystalline-]{crystalline}
\externaldocument[spaces-]{spaces}
\externaldocument[spaces-properties-]{spaces-properties}
\externaldocument[spaces-morphisms-]{spaces-morphisms}
\externaldocument[decent-spaces-]{decent-spaces}
\externaldocument[spaces-cohomology-]{spaces-cohomology}
\externaldocument[spaces-limits-]{spaces-limits}
\externaldocument[spaces-divisors-]{spaces-divisors}
\externaldocument[spaces-over-fields-]{spaces-over-fields}
\externaldocument[spaces-topologies-]{spaces-topologies}
\externaldocument[spaces-descent-]{spaces-descent}
\externaldocument[spaces-perfect-]{spaces-perfect}
\externaldocument[spaces-more-morphisms-]{spaces-more-morphisms}
\externaldocument[spaces-flat-]{spaces-flat}
\externaldocument[spaces-groupoids-]{spaces-groupoids}
\externaldocument[spaces-more-groupoids-]{spaces-more-groupoids}
\externaldocument[bootstrap-]{bootstrap}
\externaldocument[spaces-pushouts-]{spaces-pushouts}
\externaldocument[spaces-chow-]{spaces-chow}
\externaldocument[groupoids-quotients-]{groupoids-quotients}
\externaldocument[spaces-more-cohomology-]{spaces-more-cohomology}
\externaldocument[spaces-simplicial-]{spaces-simplicial}
\externaldocument[spaces-duality-]{spaces-duality}
\externaldocument[formal-spaces-]{formal-spaces}
\externaldocument[restricted-]{restricted}
\externaldocument[spaces-resolve-]{spaces-resolve}
\externaldocument[formal-defos-]{formal-defos}
\externaldocument[defos-]{defos}
\externaldocument[cotangent-]{cotangent}
\externaldocument[examples-defos-]{examples-defos}
\externaldocument[algebraic-]{algebraic}
\externaldocument[examples-stacks-]{examples-stacks}
\externaldocument[stacks-sheaves-]{stacks-sheaves}
\externaldocument[criteria-]{criteria}
\externaldocument[artin-]{artin}
\externaldocument[quot-]{quot}
\externaldocument[stacks-properties-]{stacks-properties}
\externaldocument[stacks-morphisms-]{stacks-morphisms}
\externaldocument[stacks-limits-]{stacks-limits}
\externaldocument[stacks-cohomology-]{stacks-cohomology}
\externaldocument[stacks-perfect-]{stacks-perfect}
\externaldocument[stacks-introduction-]{stacks-introduction}
\externaldocument[stacks-more-morphisms-]{stacks-more-morphisms}
\externaldocument[stacks-geometry-]{stacks-geometry}
\externaldocument[moduli-]{moduli}
\externaldocument[moduli-curves-]{moduli-curves}
\externaldocument[examples-]{examples}
\externaldocument[exercises-]{exercises}
\externaldocument[guide-]{guide}
\externaldocument[desirables-]{desirables}
\externaldocument[coding-]{coding}
\externaldocument[obsolete-]{obsolete}
\externaldocument[fdl-]{fdl}
\externaldocument[index-]{index}

% Theorem environments.
%
\theoremstyle{plain}
\newtheorem{theorem}[subsection]{Theorem}
\newtheorem{proposition}[subsection]{Proposition}
\newtheorem{lemma}[subsection]{Lemma}

\theoremstyle{definition}
\newtheorem{definition}[subsection]{Definition}
\newtheorem{example}[subsection]{Example}
\newtheorem{exercise}[subsection]{Exercise}
\newtheorem{situation}[subsection]{Situation}

\theoremstyle{remark}
\newtheorem{remark}[subsection]{Remark}
\newtheorem{remarks}[subsection]{Remarks}

\numberwithin{equation}{subsection}

% Macros
%
\def\lim{\mathop{\mathrm{lim}}\nolimits}
\def\colim{\mathop{\mathrm{colim}}\nolimits}
\def\Spec{\mathop{\mathrm{Spec}}}
\def\Hom{\mathop{\mathrm{Hom}}\nolimits}
\def\Ext{\mathop{\mathrm{Ext}}\nolimits}
\def\SheafHom{\mathop{\mathcal{H}\!\mathit{om}}\nolimits}
\def\SheafExt{\mathop{\mathcal{E}\!\mathit{xt}}\nolimits}
\def\Sch{\mathit{Sch}}
\def\Mor{\mathop{\mathrm{Mor}}\nolimits}
\def\Ob{\mathop{\mathrm{Ob}}\nolimits}
\def\Sh{\mathop{\mathit{Sh}}\nolimits}
\def\NL{\mathop{N\!L}\nolimits}
\def\CH{\mathop{\mathrm{CH}}\nolimits}
\def\proetale{{pro\text{-}\acute{e}tale}}
\def\etale{{\acute{e}tale}}
\def\QCoh{\mathit{QCoh}}
\def\Ker{\mathop{\mathrm{Ker}}}
\def\Im{\mathop{\mathrm{Im}}}
\def\Coker{\mathop{\mathrm{Coker}}}
\def\Coim{\mathop{\mathrm{Coim}}}

% Boxtimes
%
\DeclareMathSymbol{\boxtimes}{\mathbin}{AMSa}{"02}

%
% Macros for moduli stacks/spaces
%
\def\QCohstack{\mathcal{QC}\!\mathit{oh}}
\def\Cohstack{\mathcal{C}\!\mathit{oh}}
\def\Spacesstack{\mathcal{S}\!\mathit{paces}}
\def\Quotfunctor{\mathrm{Quot}}
\def\Hilbfunctor{\mathrm{Hilb}}
\def\Curvesstack{\mathcal{C}\!\mathit{urves}}
\def\Polarizedstack{\mathcal{P}\!\mathit{olarized}}
\def\Complexesstack{\mathcal{C}\!\mathit{omplexes}}
% \Pic is the operator that assigns to X its picard group, usage \Pic(X)
% \Picardstack_{X/B} denotes the Picard stack of X over B
% \Picardfunctor_{X/B} denotes the Picard functor of X over B
\def\Pic{\mathop{\mathrm{Pic}}\nolimits}
\def\Picardstack{\mathcal{P}\!\mathit{ic}}
\def\Picardfunctor{\mathrm{Pic}}
\def\Deformationcategory{\mathcal{D}\!\mathit{ef}}
\usepackage{cleveref}
\crefname{originallemma}{original lemma}{original lemmas}
\Crefname{originallemma}{Original lemma}{Original lemmas}
\makeatletter
\newlabel{topology-lemma-inverse-limit-spectral-spaces-nonempty}{{0A2W}{source}{}{sourcepointer-1}{}}
\newlabel{topology-lemma-inverse-limit-spectral-spaces-nonempty@cref}{{[originallemma][0][]0A2W}{[1][0][]source}{}{}{}}
\makeatother
\crefname{originallemma}{original lemma}{original lemmas}
\Crefname{originallemma}{Original lemma}{Original lemmas}
\makeatletter
\newlabel{topology-lemma-directed-inverse-limit-spectral-spaces}{{0A2Z}{source}{}{sourcepointer-2}{}}
\newlabel{topology-lemma-directed-inverse-limit-spectral-spaces@cref}{{[originallemma][0][]0A2Z}{[1][0][]source}{}{}{}}
\makeatother
\crefname{originallemma}{original lemma}{original lemmas}
\Crefname{originallemma}{Original lemma}{Original lemmas}
\makeatletter
\newlabel{topology-lemma-descend-opens}{{0A30}{source}{}{sourcepointer-3}{}}
\newlabel{topology-lemma-descend-opens@cref}{{[originallemma][0][]0A30}{[1][0][]source}{}{}{}}
\makeatother
\title{Sections of the pullback sheaf on the limit of a spectral system compute as the colimit of pullbacks}
\author{The Stacks Project Authors}
\date{stacks-project:sheaves.tex:local-snapshot-20261012; GNU FDL 1.2 or later}
\begin{document}
\maketitle
\setcounter{subsection}{1}
\noindent\textit{Editorial note.} All delivered mathematical text is the author's own, copied from the pinned Stacks Project chapter file sheaves.tex; no statement or proof step was rewritten, reordered or removed. The counted span is the author environment and its own proof environment. References inside the span are handed to the assembler as navigation stubs that carry the canonical tag of the cited tag only; whether each of them is load-bearing is recorded on the stub itself.\par\medskip
\noindent\textit{Allowed standard prerequisites (editorial).}\begin{itemize}
\item Topology, Lemma topology-lemma-inverse-limit-spectral-spaces-nonempty (canonical tag 0A2W): a directed inverse limit of nonempty closed subsets of spectral spaces is nonempty. cited at proof lines 3804-3806. Yields a point where the two sections differ in the limit, forcing some Z\_k empty.
\item Topology, Lemma topology-lemma-directed-inverse-limit-spectral-spaces (canonical tag 0A2Z): the limit of a directed system of spectral spaces is spectral and projections behave accordingly. cited at proof lines 3814-3816. Ensures p\_i\textasciicircum{}\{-1\}(U\_i) is quasi-compact open in the spectral space X.
\item Topology, Lemma topology-lemma-descend-opens (canonical tag 0A30): coverings of the limit descend along some transition map of a directed spectral system. cited at proof lines 3828-3830. Descends the finite covering of p\_i\textasciicircum{}\{-1\}(U\_i) to a covering of f\_a\textasciicircum{}\{-1\}(U\_i) for some a : j -> i.
\end{itemize}\medskip
\ifdefined\hypertarget\hypertarget{sourcepointer-1}{}\fi
\noindent\textit{Source reference.} Original lemma 0A2W: sources/stacks-chow/tags.txt, source lines 12719--12719. Cited at proof lines 3804-3806 for injectivity; load-bearing inside the counted span.\par\medskip
\ifdefined\hypertarget\hypertarget{sourcepointer-2}{}\fi
\noindent\textit{Source reference.} Original lemma 0A2Z: sources/stacks-chow/tags.txt, source lines 12722--12722. Cited at proof lines 3814-3816 to ensure quasi-compactness; load-bearing inside the counted span.\par\medskip
\ifdefined\hypertarget\hypertarget{sourcepointer-3}{}\fi
\noindent\textit{Source reference.} Original lemma 0A30: sources/stacks-chow/tags.txt, source lines 12723--12723. Cited at proof lines 3828-3830 to descend the covering; load-bearing inside the counted span.\par\medskip
\begin{lemma}
\label{lemma-compute-pullback-to-limit}
In the situation described above, let $i \in \Ob(\mathcal{I})$ and let
$\mathcal{G}$ be a sheaf on $X_i$. For $U_i \subset X_i$
quasi-compact open we have
$$
p_i^{-1}\mathcal{G}(p_i^{-1}(U_i)) =
\colim_{a : j \to i} f_a^{-1}\mathcal{G}(f_a^{-1}(U_i))
$$
\end{lemma}

\begin{proof}
Let us prove the canonical map
$\colim_{a : j \to i} f_a^{-1}\mathcal{G}(f_a^{-1}(U_i)) \to
p_i^{-1}\mathcal{G}(p_i^{-1}(U_i))$ is injective.
Let $s, s'$ be sections of $f_a^{-1}\mathcal{G}$
over $f_a^{-1}(U_i)$ for some $a : j \to i$. For
$b : k \to j$ let $Z_k \subset f_{a \circ b}^{-1}(U_i)$
be the closed subset of points $x$ such that the image of
$s$ and $s'$ in the stalk $(f_{a \circ b}^{-1}\mathcal{G})_x$
are different. If $Z_k$ is nonempty for all $b : k \to j$, then by
Topology, Lemma \ref{topology-lemma-inverse-limit-spectral-spaces-nonempty}
we see that $\lim_{b : k \to j} Z_k$ is nonempty too.
Then for $x \in \lim_{b : k \to j} Z_k \subset X$
(observe that $\mathcal{I}/j \to \mathcal{I}$ is initial)
we see that the image of $s$ and $s'$ in the stalk of
$p_i^{-1}\mathcal{G}$ at $x$ are different too since
$(p_i^{-1}\mathcal{G})_x = (f_{b \circ a}^{-1}\mathcal{G})_{p_k(x)}$
for all $b : k \to j$ as above. Thus if the images of $s$ and $s'$
in $p_i^{-1}\mathcal{G}(p_i^{-1}(U_i))$ are the same, then $Z_k$
is empty for some $b : k \to j$. This proves injectivity.

\medskip\noindent
Surjectivity. Let $s$ be a section of $p_i^{-1}\mathcal{G}$ over
$p_i^{-1}(U_i)$. By
Topology, Lemma \ref{topology-lemma-directed-inverse-limit-spectral-spaces}
the set $p_i^{-1}(U_i)$ is a quasi-compact open of the spectral space $X$.
By construction of the pullback sheaf, we can find an open covering
$p_i^{-1}(U_i) = \bigcup_{l \in L} W_l$,
opens $V_{l, i} \subset X_i$,
sections $s_{l, i} \in \mathcal{G}(V_{l, i})$
such that $p_i(W_l) \subset V_{l, i}$ and $p_i^{-1}s_{l, i}|_{W_l} = s|_{W_l}$.
Because $X$ and $X_i$ are spectral and $p_i^{-1}(U_i)$ is quasi-compact open,
we may assume $L$ is finite and $W_l$ and $V_{l, i}$ quasi-compact open
for all $l$. Then we can apply
Topology, Lemma \ref{topology-lemma-descend-opens}
to find $a : j \to i$ and open covering
$f_a^{-1}(U_i) = \bigcup_{l \in L} W_{l, j}$ by quasi-compact
opens whose pullback to $X$ is the covering
$p_i^{-1}(U_i) = \bigcup_{l \in L} W_l$
and such that moreover $W_{l, j} \subset f_a^{-1}(V_{l, i})$.
Write $s_{l, j}$ the restriction of the pullback of $s_{l, i}$
by $f_a$ to $W_{l, j}$. Then we see that $s_{l, j}$ and $s_{l', j}$
restrict to elements of $(f_a^{-1}\mathcal{G})(W_{l, j} \cap W_{l', j})$
which pullback to the same element $(p_i^{-1}\mathcal{G})(W_l \cap W_{l'})$,
namely, the restriction of $s$. Hence by injectivity, we can find
$b : k \to j$ such that the sections $f_b^{-1}s_{l, j}$
glue to a section over $f_{a \circ b}^{-1}(U_i)$ as desired.
\end{proof}

\end{document}
